\section{Infrastructure}
\label{sec:infrastructure}

The RSI cycle in Section~\ref{sec:rsi} links a new LoRA adapter revision to the
HCP configuration, selected data, and evaluation artifacts that produced it.
This section concerns the model-state side of that lineage and the computation
that produces the revision. \mint{} manages the model-state
lifecycle over a resident, frozen base; \longstraw{} provides a response-only
execution path when an update has a very long context; and model-dependent
controls limit rollout--training mismatch on sparse bases. These components do
not change the learned unit defined earlier: the deployable model update remains
an exported LoRA revision, paired by \mindforge{} with the HCP and evaluation
artifacts that produced it.

Figure~\ref{fig:macaron-infra-loop} separates lifecycle management from the
optional, architecture-specific execution mechanisms inside an update.
\begin{figure}[h]
\centering
\resizebox{0.98\textwidth}{!}{\begin{tikzpicture}[
    font=\sffamily\fontsize{5.8}{6.9}\selectfont,
    main/.style={draw=mindlabblue!78!black, fill=mindlabbluepale!48,
      rounded corners=2pt, line width=0.65pt, align=center,
      text width=2.18cm, minimum height=0.74cm, inner sep=4pt},
    revision/.style={main, draw=mindlabblue!90!black,
      fill=mindlabbluepale!78, line width=0.9pt},
    support/.style={draw=mindlabfg!52, fill=mindlabfg!6,
      rounded corners=1.5pt, line width=0.55pt, align=center,
      text width=3.20cm, minimum height=0.64cm, inner sep=3pt},
    band/.style={draw=mindlabfg!24, fill=mindlabfg!3,
      rounded corners=2pt, line width=0.55pt},
    flow/.style={
      -{Latex[length=1.5mm,width=1.1mm]},
      draw=mindlabfg!78, line width=0.7pt
    },
    feedback/.style={-Latex, draw=mindlabblue!62, dashed, line width=0.6pt},
    dependency/.style={-Latex, draw=mindlabfg!58, dashed, line width=0.55pt},
    label/.style={font=\sffamily\fontsize{5.3}{6.2}\selectfont,
      fill=white, inner sep=1.2pt,
      text=mindlabfg!78, align=center}
  ]
  
  \node[main] (harness) at (-6.0,0.95)
    {{\bfseries Production harness}\\HCP + tools};
  \node[main] (mindforge) at (-3.0,0.95)
    {{\bfseries MindForge}\\experience loop};
  \node[main] (mint) at (0,0.95)
    {{\bfseries MinT}\\managed update/export};
  \node[revision] (revision) at (3.0,0.95)
    {{\bfseries Adapter revision}\\versioned LoRA};
  \node[main] (serving) at (6.0,0.95)
    {{\bfseries MoL serving}\\specialist composition};
  
  \draw[flow] (harness.east) -- (mindforge.west);
  \draw[flow] (mindforge.east) -- (mint.west);
  \draw[flow] (mint.east) -- (revision.west);
  \draw[flow] (revision.east) -- (serving.west);
  
  \draw[feedback] (serving.north) -- (6.0,1.78) --
    node[label, above] {deployed interactions} (-6.0,1.78) -- (harness.north);
  
  \node[band, minimum width=7.85cm, minimum height=1.34cm]
    (supportband) at (0,-0.90) {};
  \node[anchor=west, font=\sffamily\bfseries\fontsize{5.3}{6.2}\selectfont,
    text=mindlabfg!70] at (-3.77,-0.36) {Conditional update mechanisms};
  
  \node[support] (longstraw) at (-1.95,-1.04)
    {{\bfseries LongStraw}\\response-only replay};
  \node[support] (consistency) at (1.95,-1.04)
    {{\bfseries Sparse-base controls}\\R3 + DSA alignment + filtering};
  
  \draw[dependency] (longstraw.north) -- ($(mint.south)+(-0.45,0)$);
  \draw[dependency] (consistency.north) -- ($(mint.south)+(0.45,0)$);
  \end{tikzpicture}}
\caption{\textbf{Infrastructure inside the Macaron revision loop.} This is a
conceptual artifact and control flow, not a claim that all components are
co-located or invoked for every update. \mindforge{} supplies selected
trajectories and identifies the current policy; \mint{} resolves its policy
record, manages the update, and exports an immutable LoRA adapter revision for
the \molshort{} runtime. \longstraw{} is
selected for long-context response-only updates; sparse-base consistency
controls are selected according to the model architecture. The dashed feedback
arrow denotes experience collected after deployment, not an automatic online
weight update~\citep{lu2026announcing,zhou2026longstrawlongcontextrl2m}.}
\label{fig:macaron-infra-loop}
\end{figure}
\FloatBarrier

\subsection{MinT: Model-State Lineage over a Resident Base}
\label{sec:infra_mint}

The MindLab Toolkit (\mint{})~\citep{lu2026announcing} manages LoRA RL over
resident dense and Mixture-of-Experts (MoE) base deployments. Its key
distinction is between an \emph{adapter revision}, an immutable LoRA snapshot
exported at a particular training step in serving tensor layout, and a
\emph{policy record}, the mutable service state used to resume and audit that
policy. The policy record names the compatible base version and adapter shape,
the latest trainer checkpoint and optimizer state, associated rollout records,
and the available exported revisions. Training restores mutable state on a
compatible worker, whereas rollout, evaluation, and serving select an exported
revision; trainer checkpoints and optimizer state never cross the serving
boundary directly.

This abstraction provides model-state lineage, not complete experimental
reproducibility by itself. Prompts, tools, hooks, workspace resources, task-bank
versions, evaluators, and sampling settings remain in the HCP and \mindforge{}
artifacts defined in Section~\ref{sec:rsi_mindforge}. A reported run is therefore
reconstructable at the paper's configuration boundary only when its adapter
revision and compatible base are joined with those artifacts. The narrower
\mint{} guarantee is that a rollout or score can be attributed to a fixed
adapter revision even if trainer placement and serving-cache residency change.

Adapter revisions also make the training--serving handoff smaller than a merged
checkpoint handoff. The companion \mint{} experiment compares a merge path that
materializes and loads a full checkpoint before its rollout probe with an
adapter path that loads the exported revision into a sampler whose base is
already resident. Under those stated probes, adapter-only handoff reduces the
measured handoff step by $18.3\times$ for Qwen3-4B with a rank-32 adapter and by
$2.85\times$ for Qwen3-30B with a rank-16 adapter~\citep{lu2026announcing}.
These are path-specific latency ratios, not general training-speed or inference-
throughput improvements.

The same separation supports a large addressable catalog without requiring a
large resident set. The \mint{} evaluation builds all $10^6$ entries of a packed
Qwen3-30B rank-1 adapter catalog with zero build errors, then reads a 256-entry
audit sample spanning all 100 storage shards. Serving experiments select bounded
working sets from that catalog while each actor maintains a smaller CPU cache
and GPU-active window~\citep{lu2026announcing}. Thus, $10^6$ is a measured
artifact-addressability result. It does not mean that one engine holds one
million adapters in GPU memory, nor that the entries represent one million
independently trained or behaviorally distinct policies. Cold loading remains
scheduled service work, and a revision becomes user-visible only after
compatibility checking and activation. This is the catalog boundary assumed by
the registry in Section~\ref{sec:mol_collective}~\citep{punica2024,slora2023}.

For bases that require distributed placement, \mint{} keeps model-parallel
trainer groups and inference actors resident while adapter tensors, optimizer
state, and exported revisions move between workers. The companion report
includes a Kimi K2 countdown-task LoRA RL run on a 1.04T-total/32.6B-active
parameter base using 64 H800 GPUs~\citep{lu2026announcing,kimi_k2_2025}. This is
evidence that the adapter lifecycle executes on that model and workload; it is
not evidence that \macaron{} was trained on Kimi K2, nor that the run is
reproducible from this paper alone.

\subsection{LongStraw: Response-Only Long-Context Execution}
\label{sec:infra_longstraw}

\paragraph{Response-only memory boundary.}
Long contexts make a conventional full-sequence autograd implementation
memory-bound because it retains differentiable state across the prompt and
response. For a shared prompt of length $P$ and a GRPO group of $G$ responses
with lengths $R_i$, \longstraw{}~\citep{zhou2026longstrawlongcontextrl2m}
instead captures architecture-specific prompt state without autograd and
replays one response at a time with autograd. Its live-memory boundary is
approximately
\begin{equation}
  M_{\mathrm{live}} \approx M_{\mathrm{weights}} + M_{\mathrm{prefix}}(P)
  + M_{\mathrm{grad}} + \max_i M_{\mathrm{graph}}(R_i)
  + M_{\mathrm{score}}\!\left(\sum_i R_i\right),
  \label{eq:longstraw-live-memory}
\end{equation}
The resulting graph is bounded by the longest response rather than by a
prompt-and-response graph for all group members. The full prompt still incurs
forward compute and retained model state; the method bounds graph lifetime
rather than eliminating context-dependent memory or computation.

\paragraph{Three-stage update.}
One update has three stages. First, the current policy evaluates the shared
prompt without autograd and retains only the state needed to condition response
tokens. Second, old-policy and reference scores are computed without a graph,
after which policy responses are replayed serially under autograd and each graph
is released after backward. Third, rank-local gradients are finalized once and
the optimizer steps after all $G$ members. The resulting objective is
\emph{response-only}: response tokens condition on the complete prompt, but no
gradient is propagated through prompt-token computation. ``Exact'' in the
operating points below refers to the stated token positions and this
response-only transaction, not parameter-wise equivalence to full-sequence
backpropagation.

\paragraph{Architecture-specific state.}
The retained state and distributed operators follow the base architecture. On
the hybrid recurrent/full-attention Qwen3.6-27B path, \longstraw{} retains
recurrent boundary state and context-parallel (CP) KV pages. On GLM-5.2, it
retains CP-sharded multi-head latent-attention pages and DSA indexer-key pages,
combines sparse selection across CP ranks, and dispatches routed response tokens
across expert-parallel ranks~\citep{qwen36_2026,qwen3_2025,glm52_2026,glm5_2026}.
An exact repeated update must recapture the prefix after parameters change. The
separate resident-prefix mode deliberately reuses the captured state across
optimizer steps to measure capacity and prefill amortization, so it has a
different parameter-consistency boundary.

\paragraph{Execution receipts.}
Table~\ref{tab:longstraw-support-points} condenses completed execution checks
from the \longstraw{} companion report; this paper does not rerun them as
\macaron{} experiments. They establish that the listed operations complete
under fixed eight-H20 and 32-H20 inventories, respectively. The formal 2M
receipts are Qwen at exactly $2{,}097{,}152$ total positions for $G=2$ and
$G=8$, and GLM at exactly $2{,}097{,}152$ prompt tokens for one online $G=2$
transaction. The separate Qwen $4{,}456{,}448$-position receipt is the 4.25M
resident-prefix-reuse mode. Intermediate position-ladder probes are excluded
from these reported results.

\paragraph{Evidence boundary.}
These receipts do not provide a cross-system throughput comparison, a
task-quality learning curve, or evidence that \venti{} or \tall{} was trained at
these context lengths. In particular, the Qwen rows use stored responses and
deterministic rewards, and the GLM row is one $G=2$ online DAPO-MATH transaction
with five generated tokens per response (four scored). Agent-policy quality is
evaluated separately in Section~\ref{sec:results}.

\begin{table}[H]
\centering
\caption{\textbf{Fixed-hardware \longstraw{} execution receipts, condensed from
the companion report rather than rerun in this paper.} All rows
condition response replay on a prompt captured without autograd. ``Exact''
denotes the listed token positions and completed response-only operations, not
backpropagation through prompt tokens. Here 2M denotes exactly $2{,}097{,}152$
positions or prompt tokens, and 4.25M denotes exactly $4{,}456{,}448$ positions.
Hardware and suffix workloads differ across rows, so these checks are not throughput comparisons
~\citep{zhou2026longstrawlongcontextrl2m}.}
\label{tab:longstraw-support-points}
\scriptsize
\setlength{\tabcolsep}{3.5pt}
\renewcommand{\arraystretch}{1.14}
\arrayrulecolor{mindlabfg!42}
\begin{tabularx}{\textwidth}{@{}L{0.14\textwidth}L{0.19\textwidth}L{0.24\textwidth}Y@{}}
\toprule
\rowcolor{mindlabbluepale!62}
\textbf{Check} & \textbf{Hardware / layout} & \textbf{Context / group} &
\textbf{Verified boundary} \\
\midrule
\rowcolor{mindlabbluepale!22}
\multicolumn{4}{@{}l@{}}{\textbf{Qwen3.6-27B}} \\
Exact-2M response-only replay & 8 H20\newline CP8 & 2,088,960 prompt + 8,192 response inputs\newline
= \textbf{2,097,152} positions; $G=2,8$ &
All serial response backwards and local AdamW calls complete; $dQ$ is
all-reduced. Replicated K/V-adapter synchronization is a separate audit.
Increasing $G$ from 2 to 8 adds \textbf{0.208~GB} peak allocation. \\
4.25M resident-prefix reuse & 8 H20\newline CP8 & 4,448,256 prompt + 8,192 response inputs\newline
= \textbf{4,456,448} positions; $G=8$ &
Eight optimizer steps (64 member replays) reuse one captured prefix. This is an
amortization/capacity check; the prefix is not recaptured after each update. \\
\addlinespace[2pt]
\rowcolor{mindlabbluepale!22}
\multicolumn{4}{@{}l@{}}{\textbf{GLM-5.2}} \\
Exact-2M end-to-end online & 32 H20\newline rollout TP8/PP4\newline train CP32/EP32 &
\textbf{2,097,152}-token prompt + 5 generated (4 scored) per member; $G=2$ &
One complete transaction: rollout and rewards $\rightarrow$ two 78-layer
backwards $\rightarrow$ global DSA $\rightarrow$ distributed gradient
finalization $\rightarrow$ optimizer step on all ranks. \\
\bottomrule
\end{tabularx}
\end{table}

\FloatBarrier

\subsection{Controlling Rollout--Training Mismatch on Sparse Bases}
\label{sec:infra_stability}

Separate rollout and learner engines can assign different token probabilities
even when they nominally use the same weights. Sparse MoE routing and DeepSeek
Sparse Attention (DSA) add discrete execution choices: the rollout may select
one expert set or sparse-attention index set while the learner selects another.
The resulting likelihood ratio then compares different computations, violating
the intended on-policy scoring contract~\citep{rollout_training_mismatch2025}.
The stack uses path-dependent controls to reduce this mismatch; none of them is
a general proof that rollout and learner logits are identical.

\paragraph{R3: rollout routing replay.}
For MoE paths whose rollout backend exposes expert provenance, the runtime
records the selected expert ids for every token. The learner reuses those ids
when they can be mapped to its current expert-parallel layout; if an id is
missing or unmappable, that token is excluded from the replayed policy-gradient
term rather than scored under a different expert path
~\citep{r3_moe_router2025,chiang2026routerreplay}. As a diagnostic rather than a
controlled quality ablation, the companion \mint{} report measures a mean
out-of-route scoring ratio of $0.0013\%$ with R3 across 87 logged Qwen3-30B steps,
versus $0.0097\%$ without R3 across 50 logged steps~\citep{lu2026announcing}.

\paragraph{DSA implementation alignment.}
DSA introduces a separate selection path that cannot be treated as an ordinary
linear LoRA target. The GLM integration aligns the indexer's rotary-position
layout, normalized query/key inputs, deterministic top-$k$ implementation, and
frozen-indexer defaults. It also aligns the long-context sequence- and
context-parallel layout and LoRA target-module loading before rollout
~\citep{stevenchiang2026supportglm5inmint}. These contracts remove known
implementation differences, but small numerical differences can still change a
top-$k$ index set.

\paragraph{IcePop-style residual filtering.}
The cited GLM path does not replay every DSA indexer selection. It therefore
computes the train-versus-rollout probability ratio for each token and assigns
zero importance weight when that ratio falls outside the trusted interval
recorded in the run configuration~\citep{ling_every_step2025}. This prevents a
large observed discrepancy from entering the policy-gradient term; it does not
reconstruct the rollout's sparse-attention selection or make the remaining
tokens exactly on-policy.

These controls are alternatives selected by model path, not three serial stages
applied uniformly. R3 preserves available MoE route provenance; DSA alignment
removes known implementation mismatches; and residual filtering drops detected
probability outliers when exact DSA replay is unavailable. The cited companion
studies support the mechanisms and systems diagnostics. This report does not
include a controlled \venti{} ablation that attributes benchmark gains to these
controls, so the claim here is limited to the specified execution mechanisms
and mismatch mitigation, not exact rollout--learner equivalence or downstream
quality improvement.
